% !TEX root = ../main.tex
% =====================================================================
%  §8  Derivation Of The Relation Between SW And 唐纳森 Invariants（原 PDF p. 66）
%  §9  Simple Type And Witten's Conjecture                （原 PDF p. 67）
%  公式编号 (8.1)--(8.4)、(9.1)--(9.4)，逐条保留原文编号。
% =====================================================================

\section{SW 不变量与唐纳森不变量之间关系的推导}

\begin{keytakeaways}
\item \textbf{物理图像：} 完整的唐纳森配分函数只在 $u=\infty$ 的墙上有跳变；因此 $u=\pm\Lambda^{2}$ 处库仑分支的跳变必须与希格斯分支的跳变逐墙相消，这一要求把两个分支的未知耦合锁死。
\item \textbf{数学工具：} SW 不变量自身在 $b_2^{+}=1$ 时的穿墙公式 (8.3)，与 (7.19) 给出的库仑分支跳变（一个全纯对象的留数）逐项匹配。
\item \textbf{核心结论：} 由此定出耦合 $C,P,E$ 的前导系数 (8.4)，于是 $Z_{DW}^{\xi}$ 成为只依赖 SW 不变量与经典上同调环的显式表达式，对一切 $b_2^{+}>0$ 的单连通四维流形成立。
\end{keytakeaways}

让我们回顾一下形势。我们总有
% [Tutor Note: (8.1) 的“总有”值得强调：它是按真空分支对路径积分作分解，与 $b_2^{+}$ 的取值无关。但下文整段论证只在 $b_2^{+}=1$ 时才有内容——$b_2^{+}>1$ 时没有墙，而 $Z_u^{\xi}$ 本身为零（由 (7.2) 后的讨论），(8.1) 退化为 $Z_{DW}^{\xi}=Z_{IR,Higgs}$，这正是第 9 节 (9.1) 成立的直接原因。]
\begin{equation}
Z_{DW}^{\xi}(p,\Sigma)=Z_u^{\xi}(p,\Sigma)+Z_{IR,Higgs}
\tag{8.1}
\end{equation}
当 $X$ 有 $b_2^{+}=1$ 时，我们知道 $Z_u^{\xi}(p,\Sigma)$ 作为 $\omega\in H^{2}(X;\RR)$ 的函数，会在来自 $u=\infty$ 与 $u=\pm\Lambda^{2}$ 奇异性的各道墙 $W(\lambda)$ 上有间断。

此外，$Z_{DW}^{\xi}(p,\Sigma)$ 也有间断，而这些间断完全可由 $Z_u^{\xi}(p,\Sigma)$ 来自 $u=\infty$ 的间断解释。

因此，跨过所有墙 $W(\lambda)$ 必须有
% [Tutor Note: (8.2) 的逻辑是“真正的配分函数之跳变必须由 $u=\infty$ 的库仑分支跳变完全解释”，故 $u=\pm\Lambda^{2}$ 处库仑分支与希格斯分支的跳变必须逐墙相消。注意这是对每一道墙、每一个 $\lambda$ 分别成立的等式，强于“总和不变”；它的实际用途是：左边由 (7.19) 的留数显式已知，右边由 (6.7) 的一般形式给出，于是可以反解出 $C,P,E$。]
\begin{equation}
0=\Delta_{u=\Lambda^{2},\lambda}Z_{IR,Coul}+\Delta_{u=\Lambda^{2},\lambda}Z_{IR,Higgs}
\tag{8.2}
\end{equation}
的确，从数学上已知：当 $X$ 有 $b_2^{+}=1$ 时，SW 不变量 $SW(\lambda)$ 不再是（与 $\omega$ 无关的）不变量，它会在由自旋$^{c}$结构决定的各道墙 $W(\lambda)$ 上改变。其穿墙公式（WCF）特别简单：
% [Tutor Note: (8.3) 两端是 $\omega\cdot\lambda=0$ 两侧的极限，跳变值为 $(-1)^{1+n(\lambda)}$；按墙的定向与跨墙方向，这一跳变也可写成 $\pm(-1)^{1+n(\lambda)}$，原文取定“正侧减负侧”从而把符号写死为 $(-1)^{1+n(\lambda)}$。又因 $n(\lambda)$ 由 (6.3) 定义、且 $\lambda\to-\lambda$ 不改变它，该跳变在自旋$^{c}$结构整体反号下不变，这是它能成为良好穿墙公式的前提。]
\begin{equation}
SW(\lambda)\big|_{\omega\cdot\lambda=0^{+}}
-
SW(\lambda)\big|_{\omega\cdot\lambda=0^{-}}
=
(-1)^{1+n(\lambda)}
\tag{8.3}
\end{equation}
从数学上看，在这样的墙上存在 SW 方程的解且 $M=0$：这是一个在整体 $U(1)$ 规范变换下不变的可约解，模空间因而变奇异。从物理上看，在这些墙上由于 $M=0$，希格斯分支与库仑分支可以“混合”。

现在我们可以计算 $\Delta_{u=\Lambda^{2},\lambda}Z_{IR,Coul}$，因为我们有它的显式表达式 (7.19)；而已知希格斯分支贡献的一般形式 (6.7) 与 SW 的穿墙公式 (8.3)，我们就能算出未知耦合 $C(u),P(u),E(u)$。例如我们得到
\begin{equation}
\begin{aligned}
C(u) &= \left(\frac{a_d}{q_d}\right)^{1/2}=4\ee^{\ii\pi/4}+\OO(a_d)\\
P(u) &= \ee^{\ii\pi/32}2^{5/4}+\OO(a_d)\\
E(u) &= \ee^{\ii\pi/8}2^{3/4}+\OO(a_d)
\end{aligned}
\tag{8.4}
\end{equation}
其中 $q_d=\ee^{2\pi\ii\tau_d}$。这些都是由模函数决定的、完全显式已知的级数。
% [Tutor Note: 这三个常数不能由 (6.7) 或 (7.19) 单独定出，而要用“两边在同一道墙上的跳变必须相同”逐项匹配：把 (7.19) 在各墙处的留数（已知的模形式数据，见 (7.21) 与 (7.24)）与 (6.7) 中 $\mathrm{Res}_{a_d=0}$ 的表达式相等，即可解出 $C,P,E$ 在 $u=\Lambda^{2}$ 附近的 $a_d$ 展开系数。原文只列前两项，是因为后续各节只用到前导项。注意 $C(u)=(a_d/q_d)^{1/2}$ 是平方根型：它在 $u=\Lambda^{2}$（$a_d=0$，同时 $q_d\to0$）处取有限非零值 $4\ee^{\ii\pi/4}$，这说明 $a_d$ 与 $q_d$ 在该点以相同的速率趋于零；这也解释了 (6.2) 中 $C(u)^{\lambda^{2}}$ 如何给出 $\lambda^{2}$ 次的整幂次交截数。]

\emph{总结如下：}我们现在有了 $Z_{DW}^{\xi}(p,\Sigma)$ 的一个完全显式的表达式，它由 SW 不变量与经典上同调环表出。该表达式对一切 $b_2^{+}>0$ 的单连通四维流形成立，并且可以容易地推广到非单连通的情形。\footnote{〔原稿边注〕距第 6 节已过 $51$ 分钟。}

\section{单型与威滕猜想}

\begin{keytakeaways}
\item \textbf{物理图像：} SW 单型（只有 $n(\lambda)=0$ 的 $\lambda$ 才使模空间非空）意味着 $\cM(\lambda)$ 只是有限个定向点，$Z_{IR,Coul}=0$，全部贡献来自 $u=\pm\Lambda^{2}$ 两个阿贝尔真空，即磁单极点与双荷子点。
\item \textbf{数学工具：} $n(\lambda)=0$ 时 (6.7) 的留数直接给出 $SW(\lambda)$ 本身，再配合 (8.4) 定出的耦合，即得只含 SW 不变量与经典上同调环的闭式；此外是克龙海默--姆罗夫卡 (克龙海默--姆罗夫卡) 的微分算子 $\partial^{2}/\partial p^{2}-4$。
\item \textbf{核心结论：} 威滕猜想 (9.1)（整体指数 $2^{\frac{1}{4}(7\chi+11\sigma)}$，两个求和分别来自磁单极与双荷子点）给出配分函数的闭式；并由此推出“SW 单型 $\Rightarrow$ KM 单型”以及广义单型的阶 $r=1+\max_{\lambda}n(\lambda)$。
\end{keytakeaways}

关于四维流形 $X$ 上塞伯格--威滕不变量的一个关键性质是：$\cM(\lambda)$ 只对有限多个 $\lambda$ 非空。这可由 Weitzenböck 型论证得出：取塞伯格--威滕方程的平方和即可。
% [Tutor Note: 该论证的要点是把 SW 方程（含 $D_A M=0$ 与 $F^{+}=\sigma(M)$ 型方程）平方后积分，得到形如 $\int(|D_AM|^{2}+|F^{+}-\cdots|^{2}+s|\psi|^{2})=0$ 的恒等式；一旦某些曲率项占优（例如 $b_2^{+}>1$ 时对 $\omega$ 取适当组合），就强迫 $M=0$，从而把可能的 $\lambda$ 限制在有限集。原文只说“Weitzenbock 型”，未给细节，属讲义惯例。]

现在定义：若 $\cM(\lambda)\neq\varnothing$ 只对满足 $n(\lambda)=0$ 的 $\lambda$ 成立，则称 $X$ 为 \emph{塞伯格--威滕单型}。此时 $\cM(\lambda)$ 是有限个定向点的并。计算 $SW(\lambda)$ 时，我们就是在字面意义上数方程解的个数，一如本讲开头所做的那样。

对 $b_2^{+}>1$，SW 模空间依赖度规，但其配边型不变。于是 SW 不变量 $SW(\lambda)$ 与度规无关。\emph{塞伯格--威滕基本类}是使 $SW(\lambda)\neq0$ 的那些自旋$^{c}$结构 $\lambda$。
% [Tutor Note: 〔原稿边注〕原文此处的边注是：“This is not the same as the set of $\lambda$ for which $\cM(\lambda)\neq\varnothing$. Clarify.”（这与使 $\cM(\lambda)\neq\varnothing$ 的 $\lambda$ 集合并不相同，待澄清。）作者提醒：基本类是 $SW(\lambda)\neq0$ 的类，而“模空间非空”只要求 $SW$ 的某个交截数不显然为零，两者不等价（例如 $\cM(\lambda)$ 非空但 $n(\lambda)>0$ 时，插入次数不匹配可使 $SW(\lambda)=0$）。]

一个奇怪的事实是：所有已知的 $b_2^{+}>1$ 的单连通 $X$ 都是塞伯格--威滕单型的，但没有证明表明所有这样的 $X$ 必为单型。

无论如何，现在假定 $X$ 是 SW 单型，并且 $b_2^{+}>1$。此时 $Z_{IR,Coul}=0$，积分完全由 $u=\pm\Lambda^{2}$ 处的贡献给出。而且由于 $n(\lambda)=0$，这些贡献容易算出。综合起来，我们就得到 \cite{ref55} 的关键论断，数学文献中称之为“威滕猜想”：
% [Tutor Note: 为什么 $b_2^{+}>1$ 时 $Z_{IR,Coul}=0$：由 (7.2) 的讨论，$\eta$ 只有一个零模、$\chi$ 有 $b_2^{+}$ 个零模，而作用量中的 $\eta\chi$ 双线性只能吃进一个 $\chi$ 零模；$b_2^{+}>1$ 时格拉斯曼积分恒为零。于是 (8.1) 退化为 $Z_{DW}^{\xi}=Z_{IR,Higgs}$，而希格斯分支的贡献只来自两个阿贝尔真空——这正是 (9.1) 只含两个求和的原因。]
\begin{equation}
\begin{aligned}
Z_{DW}^{\xi}(p,\Sigma)
={}&
2^{\frac{1}{4}(7\chi+11\sigma)}
\Bigg(
e^{\frac{1}{2}\Sigma^{2}+2p}
\sum_{\lambda\in\Gamma_w}SW(\lambda)\,e^{2\pi\ii(\lambda\cdot\lambda_0+\lambda_0^{2})}\,e^{2\Sigma\cdot\lambda}\\
&+
\ii^{\chi_h}\,
e^{-\frac{1}{2}\Sigma^{2}-2p}
\sum_{\lambda\in\Gamma_w}SW(\lambda)\,e^{2\pi\ii(\lambda\cdot\lambda_0+\lambda_0^{2})}\,e^{-\ii2\Sigma\cdot\lambda}
\Bigg)
\end{aligned}
\tag{9.1}
\end{equation}
其中 $\chi_h:=(\chi+\sigma)/4$。
% [Tutor Note: 指数 $\frac14(7\chi+11\sigma)$ 的来历：它汇总了 $\alpha^{\chi}\beta^{\sigma}$ 与两点处 $\Delta^{\sigma/8}(\dd u/\dd a)^{\chi/2}$ 及留数给出的幂次。由 $\chi+\sigma=4$ 可把它改写为 $7+\sigma$，即整体归一化因子只由 $\chi$ 与 $\sigma$ 的组合固定；这一步与 (7.24)、(8.4) 的匹配是同一次计算的两个侧面。另有 $\chi_h=(\chi+\sigma)/4\in\ZZ_{\ge0}$（因 $\chi+\sigma\equiv0\bmod4$），故 $\ii^{\chi_h}$ 是 $4$ 次单位根，第二项的相位良定义。]

这里我们取 $\Lambda=1$。第一个求和来自磁单极点 $u=\Lambda^{2}$，第二个求和来自双荷子点 $u=-\Lambda^{2}$。注意我们写出的是\emph{物理}配分函数。克龙海默--姆罗夫卡的生成函数比它大一个因子 $2$。
% [Tutor Note: 两个求和分别来自磁单极点与双荷子点：在 $u=\Lambda^{2}$ 处 $a_d=0$、磁单极不带电，贡献给出 $e^{+\frac12\Sigma^{2}+2p}$ 一支；在 $u=-\Lambda^{2}$ 处（经第 5.4.2 节的 $\ZZ_2$ 对称性 $u\to-u$、$a\to\ii a$ 联系）参数换成对偶变量，指数变号、$\Sigma\cdot\lambda$ 也带因子 $\ii$，故第二支为 $e^{-\frac12\Sigma^{2}-2p}e^{-\ii2\Sigma\cdot\lambda}$，并带相对相位 $\ii^{\chi_h}$。与克龙海默--姆罗夫卡生成函数相差因子 $2$，源于他们把两个点同等计数而物理配分函数只计一次。]

现在，克龙海默与姆罗夫卡在分析唐纳森多项式时引入了\emph{单型}的概念——我们称之为 KM 单型。它说的是配分函数 $Z_{DW}$ 满足如下简单微分方程：
\begin{equation}
\left(\frac{\partial^{2}}{\partial p^{2}}-4\right)Z_{DW}^{\xi}(p,\Sigma)=0
\tag{9.2}
\end{equation}
我们注意到，由我们的一般物理表达式，立即可得：对 $b_2^{+}>1$，若 $X$ 是 SW 单型，则它必为 KM 单型。
% [Tutor Note: $(\partial_p^{2}-4)$ 作用在 $\ee^{\pm2p}$ 上为零：$(\partial_p^{2}-4)\ee^{2p}=0$。而 (9.1) 中每一项对 $p$ 的依赖恰是 $\ee^{2p}$ 或 $\ee^{-2p}$ 乘以不依赖 $p$ 的常数（SW 单型时 $n(\lambda)=0$，模空间积分不产生 $p$ 的多项式因子），故 (9.2) 逐项成立。这就是“SW 单型 $\Rightarrow$ KM 单型”的全部内容——注意这不是数学上的分析结果，而是 (9.1) 的直接推论。]

KM 还引入了\emph{广义单型}的概念。它说的是：对某个 $r$，
\begin{equation}
\left(\frac{\partial^{2}}{\partial p^{2}}-4\right)^{r}Z_{DW}^{\xi}(p,\Sigma)=0
\tag{9.3}
\end{equation}
注意，我们有一个物理证明：所有 $b_2^{+}>1$ 的单连通四维流形都是广义 KM 单型的。这是“基本类只有有限多个”这一事实的简单推论。因此我们可以取
\begin{equation}
r=1+\max_{\lambda}n(\lambda)
\tag{9.4}
\end{equation}
% [Tutor Note: $r=1+\max_\lambda n(\lambda)$ 的机制：对形如 $\ee^{2p}f(p)$ 的项，$(\partial_p^{2}-4)\left(\ee^{2p}f\right)=\ee^{2p}\left(f''+4f'\right)$，每作用一次就把 $f$ 的次数降低 $1$；而对 $n(\lambda)>0$ 的类，$\cM(\lambda)$ 上的交截数给出 $p$ 的至多 $n(\lambda)$ 次的多项式因子。故作用 $1+\max_\lambda n(\lambda)$ 次后所有项归零。这里多出的“$1$”来自每一项必带的整体因子 $\ee^{\pm2p}$。基本类有限保证 $\max_\lambda n(\lambda)$ 有限，这正是 (9.4) 能取到有限 $r$ 的原因。]

% [Tutor Note: 【补遗】原稿第 66 页末（(9.4) 之后、§10 之前）尚有一段以 "Remark" 起首
%  的未编号文字，内容是关于 ZDW 局部化到 N=1 真空的解释，作者标注 ``ADD MORE ABOUT THIS.''
%  尚未写完。为不使译文脱漏原文内容，此处照译补入；它不属于任何编号公式。]

\begin{remark*}
\textbf{注记。} 关于 $Z_{\mathrm{DW}}$ 的局域化，有一个漂亮的解释，即把它理解为向
$\mathcal{N}=1$ 真空的局域化。

基本想法是：有时（例如在凯勒流形上）可以给费米子加上一个质量项，使
$\mathcal{N}=2$ 破缺到 $\mathcal{N}=1$，但仍保持一个拓扑对称性。参见 \cite{ref54}。
\footnote{〔原稿边注〕作者在此处旁注 ``ADD MORE ABOUT THIS.''（此处应再多写一些），
说明本段尚未展开。}
\end{remark*}
